\section{Los bioelementos de la materia viva}

La vida se compone de ciertos bloques de construcción específicos, llamados biomoléculas

\section{El agua y la vida}

El agua es otro elemento muy importante para la vida. La unidad de la vida, la célula, se compone de aproximadamente 80\% agua, por lo que es la molécula más abundante en los seres vivos. Esta abundancia no surge de la nada, y es que el agua es tan abundante en los seres vivos porque sus características fisicoquímicas ejercen una profunda influencia en la estructura, organización y funcionamiento de estos seres vivos.   

\begin{figure}[H]
    \centering
    \vspace{6pt}
    \LARGE \chemfig[fixed length=true]{
        \Charge{150[anchor=135+\chargeangle]=\small$\delta^{+}$}{H}
        -[:37.75]
        \Charge{90[anchor=180+\chargeangle]=\small$\delta^{-}$}{O}
        -[:-37.75]
        \Charge{30[anchor=180+\chargeangle]=\small$\delta^{+}$}{H}}
    \caption{Representación sencilla de una molécula de agua}
    \label{fig:molécula_agua}
\end{figure}

\definition{Electronegatividad}{Es la capacidad de un átomo o grupos funcionales para atraer electrones}

La geometría de la molécula de agua, al igual que los valores de electronegatividad del hidrógeno ($En=2.2$) y del oxígeno ($En=3.44$), favorece que se creen cargas parciales alrededor de los átomos que la conforman, adquiriendo los átomos de hidrógeno una carga parcial positiva ($\delta +$) y el átomo de oxígeno una carga pacial negativa ($\delta -$).

Estas cargas parciales dan origen a uno de los enlaces más importantes en la naturaleza: los puentes de hidrógeno.

Estos enlaces se forman cuando la parte parcialmente positiva de los átomos de hidrógeno de una molécula de agua se orientan hacia la parte parcialmente negativa de un átomo de oxígeno de otra molécula de agua. Si bien este tipo de enlaces pueden aparecer en otros compuestos donde el átomo con carga parcial negativa es nitrógeno, así como en algunas biomoléculas (como se verá más adelante), el agua es la molécula con mayor capacidad de formar este tipo de enlaces.

Es a partir de estos puentes de hidrógeno que podemos comprender muchas de las propiedades fisicoquímicas del agua y, por ende, el porqué es tan importante para la vida. 

El agua líquida a 0 °C tiene 15% menos puentes de hidrógeno que el hielo; también se
sabe que los puentes de hidrógeno en el agua líquida están en un proceso continuo de
rompimiento y restablecimiento. Se trata de un sistema dinámico, con rápidas
fluctuaciones, en donde las moléculas individuales tienen participaciones cambiantes. El
modelo que mejor resume y adapta las propiedades del agua incluye poco menos de 80%
de moléculas que tiene la estructura del hielo (figura 2-3); 20% de las moléculas se
empacan más y son más densas que las del hielo. El mayor empaque de las moléculas de
agua se logra al ocupar de manera parcial los “huecos” que quedan, según se descubre al
analizar de manera cuidadosa la estructura del hielo. Lo anterior explica por qué el hielo
tiene densidad menor (0.92 g/mL) que el agua líquida a 0 °C (1.00 g/mL). El hecho de
que el hielo flote en los océanos, debido a su menor densidad, favorece en grado notable
el calentamiento de enormes masas de agua por el Sol, con lo que la temperatura de la
Tierra se mantiene más alta y con menores fluctuaciones.


Figura 2-3. Estructura del hielo. Se observa el arreglo tetraédrico de las moléculas de agua debido a la existencia
de los puentes de hidrógeno entre las propias moléculas. Es claro que el arreglo tetraédrico no es rígido, lo cual
permite lo que se ha llamado una estructura “abierta” del agua en el hielo, que explica la menor densidad del hielo en relación con la del agua líquida.

\subsection{Propiedades fisicoquímicas del agua líquida}

Las propiedades del agua están asociadas y determinadas de manera estrecha por sus puentes de hidrógeno intermoleculares. A continuación se revisan de forma breve algunas de sus propiedades fisicoquímicas, incluyendo ejemplos de la repercusión que dichas propiedades tienen en biología.

\begin{itemize}
    \item Calor específico
    \item Calor de fusión
    \item Calor de evaporación
    \item Tensión superficial y adhesión 
    \item Constante dieléctrica
    \item Hidratación
\end{itemize}

\subsubsection{Calor específico}

El calor específico se define como la cantidad de energía calorífica necesaria para
aumentar la temperatura de 1 g de una sustancia en 1°C. Es alto para el agua (1 cal/g) en
comparación con otros líquidos. Esta propiedad se puede entender como que el agua,
gramo por gramo, absorbe más energía calorífica por grado que la mayoría de las demás
sustancias. Se aprovecha esta propiedad al usar el agua como enfriador en los motores de
automóviles y en los sistemas de calefacción de los edificios. La humedad de los bosques
es definitiva para mantener con menores cambios de temperatura estos ecosistemas, en
comparación con lo que se observa en los desiertos. En los mamíferos, el elevado calor
específico del agua ayuda a mantener la temperatura homogénea en el cuerpo, mediante
el bombeo constante de sangre desde el corazón hacia los tejidos periféricos, debido a
que el componente más abundante de la sangre es el agua.

\subsubsection{Calor de fusión}

Se llama calor molar de fusión a la energía empleada en la fusión de un mol de un sólido.
Se mide en el punto de fusión del sólido, y para el agua esto ocurre a 0 °C. De nuevo es
posible comprobar que, en el paso de hielo a agua, el valor resulta de forma comparativa
alto (80 cal/g). El calor de fusión representa la energía cinética que deben adquirir las
moléculas del sólido para pasar de un orden continuo, impuesto por las fuerzas de
atracción en el sólido, hacia un orden discontinuo, característico del líquido. En los seres
vivos, el alto calor de fusión del agua ofrece un sistema eficiente de protección contra el
congelamiento. Para congelar un mol de agua (18 g) se necesita eliminar la misma
cantidad de energía que se absorbe al descongelarlo.

\subsubsection{Calor de evaporación}

Se llama calor molar de evaporación a la energía que se invierte en la evaporación de un
mol de un líquido en su punto de evaporación. Como en el caso del calor de fusión, el de
evaporación representa la cantidad de energía que requieren las moléculas en el estado
líquido para vencer su mutua atracción y alejarse unas de otras, como en los gases. El
calor de evaporación del agua también es alto; por lo tanto, se minimizan las pérdidas de
agua que pudieran ocurrir en los seres vivos debido a la evaporación, de manera que
protege contra la deshidratación. Además, si ocurre la evaporación, provee de un
eficiente sistema de enfriamiento, debido a que la energía indispensable para la

77

evaporación “la toma” el agua de la superficie del ser vivo, con lo que se aprecia una
sensación de frescura. Ésta es una situación que se presenta, por ejemplo, en un clima
cálido, en el que ocurre sudoración profusa con evaporación del agua presente en el
sudor.

\subsubsection{Tensión superficial y adhesión}

La tensión superficial se manifiesta en la superficie de un líquido y corresponde a la
cantidad de energía necesaria para aumentar su superficie por unidad de área. La tensión
superficial depende de la atracción que sufren las moléculas de la superficie hacia el seno
del líquido. Las fuerzas de atracción ocurren entre todas las moléculas; así, las moléculas
ubicadas en el seno del líquido son atraídas de manera mutua en todas direcciones,
mientras que las moléculas de la superficie, al no existir moléculas del líquido que las
atraigan hacia “afuera”, no tienen esa atracción de compensación. Por lo tanto,
predomina, con mucho, la atracción hacia el seno del líquido. El resultado es la tendencia
a formar gotas en las que se exhiba la menor superficie compatible con la fuerza de
gravedad. La tensión superficial relativamente alta del agua se encuentra influida por su
adhesividad y por su viscosidad. Por adhesividad se entiende la fuerza de unión con una
superficie, y por viscosidad, la resistencia a fluir a través de un tubo capilar. Un líquido
es menos viscoso cuanto más rápido fluya. La viscosidad del agua es alta en relación con
su peso molecular.


La combinación de estas tres propiedades del agua (tensión superficial, adhesividad y
viscosidad) al parecer tiene importancia especial en los seres vivos. Se nota en el caso de
los capilares. En un capilar, la fuerza adhesiva hace que el agua suba por él y, debido a la
tensión superficial y viscosidad de la misma agua, las moléculas del seno del líquido
“siguen” a las de la superficie en su ascenso por el capilar. El proceso no depende de una
fuente externa de energía y su límite está dado por el diámetro del capilar y la fuerza de
gravedad; este fenómeno se encuentra muy difundido en la naturaleza, por ejemplo, en
las raíces y en los tallos de las plantas. Es notable cómo el agua puede ascender más de
100 m de altura en los árboles, desde abajo del nivel del piso hasta la copa del árbol, sin
requerir de un sistema de bombeo.

\subsubsection{Constante dieléctrica}

Es la propiedad de los solventes de separar iones de cargas
opuestas. Al igual que los valores de las propiedades fisicoquímicas descritas arriba, la
constante dieléctrica del agua a 20

oC es muy grande, de 80.

\subsubsection{Hidratación}

Es la capacidad de rodear los iones con moléculas de agua que se
orientan según la carga de los iones y se disponen en capas concéntricas de moléculas
alrededor del ion. Cuando el solvente es distinto del agua, esta propiedad recibe el
nombre de solvatación.

\subsection{El agua como solvente}

Al agua se le identifica como el solvente universal, debido a su capacidad de disolver más
sustancias y en mayor cantidad que cualquier otro solvente. La constante dieléctrica del
agua, su capacidad de hidratación y de formación de puentes de hidrógeno, así como la
posibilidad de romper los enlaces iónicos de las moléculas que disuelve, explican el
comportamiento del agua como solvente universal.
La atracción electrostática de iones con cargas opuestas disminuye 80 veces (constante
dieléctrica del agua) al colocarse en agua. Un ion positivo en el seno del agua atrae las
cargas negativas de las moléculas del agua, las cuales se disponen en capas alrededor del
ion. Si la carga del ion inmerso es negativa, atraerá las cargas positivas de las moléculas
de agua. El resultado es que los iones se separan y quedan disueltos (figura 2-4).

Figura 2-4. Esquema de la solvatación de iones por moléculas de agua que adquieren una orientación definida.
Cuando las moléculas que se colocan en el agua contienen puentes de hidrógeno entre
ellas, las moléculas de agua tienden a debilitar tales puentes y sustituirlos por nuevos
puentes de hidrógeno, en los que intervienen moléculas de agua; con ello, se separan
entre sí las moléculas del compuesto que se colocó en el agua. Un ejemplo sería la
disolución de la sacarosa (azúcar de caña) en agua.
Además, el agua manifiesta una importante influencia en las moléculas que no se
disuelven en ella, las no polares del tipo de las grasas o aceites. Si varias gotitas de grasa
se colocan en el agua, cada una de ellas ocupará un espacio que promoverá un
incremento en la organización del agua alrededor de la superficie hidrofóbica, con la
consecuente disminución de la entropía del sistema. El resultado final, compatible con la
situación más estable, es que se formen dos fases, la de arriba, que contiene a las
moléculas no polares, y la de abajo, formada por las moléculas de agua. De esta manera,

80

hay una gran disminución de la superficie de contacto entre el agua y la región
hidrofóbica, lo que conduce a un aumento de la entropía. A este fenómeno se le conoce
como efecto hidrofóbico, y se basa en la tendencia del agua a rechazar las superficies no
polares. Es decir, aunque existen fuerzas de atracción entre las moléculas no polares
(fuerzas de van der Waals, fuerzas de dispersión de London), la “fuerza” que dirige a las
moléculas no polares a separarse del agua, es el rechazo del agua por las superficies
hidrofóbicas. El acomodo de los lípidos para formar membranas en las células o el
ocultamiento de los residuos de aminoácidos en el núcleo de una proteína siguen las
tendencias generales aquí esbozadas.
En resumen, la capacidad del agua de afectar la conformación de todas las moléculas
celulares, unas por mantenerlas en la solución y otras por eliminarlas de la solución, hace
del agua el medio idóneo para que todas las moléculas presentes en los seres vivos
adopten las formas que tienen: las más estables y las que les permiten realizar sus
funciones en los seres vivos.